\documentclass[12pt,a4paper]{article}

\usepackage[utf8]{inputenc}
\usepackage[T1]{fontenc}
\usepackage[spanish]{babel}
\usepackage{amsmath, amssymb}
\usepackage{geometry}
\usepackage{booktabs}
\usepackage{array}
\usepackage{xcolor}
\usepackage{hyperref}
\usepackage{enumitem}
\usepackage{fancyhdr}
\usepackage{graphicx}
\usepackage{listings}
\usepackage{tcolorbox}
\usepackage{multirow}
\usepackage{caption}

\geometry{top=2.5cm, bottom=2.5cm, left=3cm, right=3cm}

% ── Colores ──────────────────────────────────────────────────────────────────
\definecolor{uddblue}{RGB}{0, 75, 135}
\definecolor{lightgray}{RGB}{245, 245, 245}
\definecolor{codegray}{RGB}{100, 100, 100}
\definecolor{codegreen}{RGB}{0, 120, 0}
\definecolor{codeblue}{RGB}{0, 0, 180}
\definecolor{codered}{RGB}{160, 0, 0}

% ── Estilo de código Python ───────────────────────────────────────────────────
\lstset{
    language     = Python,
    backgroundcolor = \color{lightgray},
    basicstyle   = \small\ttfamily,
    keywordstyle = \color{codeblue}\bfseries,
    commentstyle = \color{codegreen}\itshape,
    stringstyle  = \color{codered},
    numberstyle  = \tiny\color{codegray},
    numbers      = left,
    stepnumber   = 1,
    breaklines   = true,
    frame        = single,
    rulecolor    = \color{gray!40},
    tabsize      = 4,
    showstringspaces = false,
    literate =
        {á}{{\'{a}}}1 {é}{{\'{e}}}1 {í}{{\'{i}}}1
        {ó}{{\'{o}}}1 {ú}{{\'{u}}}1 {ñ}{{\~{n}}}1
        {Á}{{\'{A}}}1 {É}{{\'{E}}}1 {Í}{{\'{I}}}1
        {Ó}{{\'{O}}}1 {Ú}{{\'{U}}}1 {Ñ}{{\~{N}}}1
}

% ── Cajas de ayuda ────────────────────────────────────────────────────────────
\tcbuselibrary{skins, breakable}
\newtcolorbox{ayuda}[1]{
    colback   = lightgray,
    colframe  = uddblue,
    fonttitle = \bfseries\small,
    title     = {#1},
    breakable,
    left      = 6pt,
    right     = 6pt,
    top       = 4pt,
    bottom    = 4pt
}

% ── Encabezado y pie ─────────────────────────────────────────────────────────
\pagestyle{fancy}
\fancyhf{}
\fancyhead[L]{\small\color{gray} IIP314W\,--\,Optimización Aplicada a Negocios}
\fancyhead[R]{\small\color{gray} Tarea 2\,--\,2026-T2}
\fancyfoot[C]{\small\thepage}
\renewcommand{\headrulewidth}{0.4pt}

\hypersetup{colorlinks=true, linkcolor=uddblue, urlcolor=uddblue}

% ── Espaciado entre párrafos ──────────────────────────────────────────────────
\setlength{\parskip}{0.5em}
\setlength{\parindent}{0pt}

% =============================================================================
\begin{document}
% =============================================================================

% ── Título ───────────────────────────────────────────────────────────────────
\begin{center}
    {\Large\bfseries\color{uddblue}
     (IIP314W) Optimización Aplicada a Negocios}\\[0.4em]
    {\large Universidad del Desarrollo}\\[0.8em]
    \noindent\rule{\linewidth}{1.5pt}\\[0.5em]
    {\LARGE\bfseries Tarea 2}\\[0.3em]
    {\large Modelamiento y Optimización:\\
     Localización de Servidores en la Nube (Netflix\,--\,AWS)}\\[0.5em]
    \noindent\rule{\linewidth}{1.5pt}\\[0.8em]
    \begin{tabular}{@{}ll@{}}
        \textbf{Profesor:} & Ing.\ Rodrigo Trigo Vilches \\[0.3em]
        \textbf{Ayudante:} & Lic.\ Vicente Ramírez Almonacid \\[0.3em]
        \textbf{Período:}  & Año 2026\,--\,Trimestre 2
    \end{tabular}
\end{center}

\vspace{1em}
\noindent\rule{\linewidth}{0.4pt}
\vspace{0.5em}

% =============================================================================
\section*{Contexto: Servicios en la Nube para Netflix}
% =============================================================================

Usted prestará servicios informáticos a la compañía de \emph{streaming}
\textbf{Netflix}. Para ello alojará servidores en los centros de datos de
\textbf{AWS} (Amazon Web Services) y debe decidir \textbf{dónde y cuántos
servidores dar de alta}, de modo de \textbf{minimizar el costo total} y, al
mismo tiempo, cumplir con los requerimientos de negocio (atender a todos los
clientes, respetar capacidades y garantizar una calidad de servicio mínima).

Netflix presta servicios en $n$ países. En cada país $i$ existe una masa de
$C_i$ clientes que deben ser atendidos. AWS cuenta con $z$ \textbf{zonas de
servicio} (cada zona representa un país). Como es habitual, $z < n$: solo los
primeros $z$ países poseen una zona AWS local, mientras que los $(n-z)$ países
restantes \textbf{no tienen infraestructura propia} y deben ser atendidos desde
otra zona. Para simplificar la notación, los países con zona AWS se ubican
\textbf{al comienzo} de la lista y los países sin zona, al final.

Dar de alta un servidor en la zona $j$ tiene un costo $S_j$ y puede atender
hasta $U_j$ clientes a su máxima capacidad. En cada zona se puede instalar un
máximo de $M_j$ servidores. Adicionalmente, \textbf{utilizar} una zona (es
decir, operar al menos un servidor en ella) implica un \textbf{costo
administrativo fijo} $A_j$ (licencias, personal de soporte, cumplimiento
regulatorio): si una zona no se utiliza, este costo se ahorra por completo.

Existe además una restricción de \textbf{latencia}: por calidad de servicio,
los clientes de un país solo pueden ser atendidos desde zonas
\textbf{geográficamente compatibles} (una región no puede servirse desde otra
sin degradar la experiencia). Finalmente, los países que \emph{sí} poseen zona
local están sujetos a un \textbf{acuerdo de nivel de servicio (SLA)}: al menos
una fracción $\beta$ de sus clientes debe ser atendida localmente.

\begin{ayuda}{Trabajo en grupo}
\vspace{2pt}
La tarea se desarrolla \textbf{en parejas o grupos de 3} (ni más ni menos).
Tanto el informe como el código deben incluir los nombres de todos los
integrantes. \textbf{Es obligatorio resolver el modelo con \texttt{gurobipy}}.
\end{ayuda}

\noindent\rule{\linewidth}{0.4pt}

% =============================================================================
\section*{Datos del problema}
% =============================================================================

Para esta tarea, Netflix opera en $n = 8$ países, de los cuales $z = 3$ poseen
zona AWS local. Todas las cantidades de clientes y capacidades están expresadas
\textbf{en miles de clientes}; los costos, en \textbf{miles de USD mensuales}.

\begin{center}
\renewcommand{\arraystretch}{1.2}
\begin{tabular}{clcc}
\toprule
\textbf{$i$} & \textbf{País} & \textbf{¿Zona AWS local?} &
\textbf{Clientes $C_i$ (miles)} \\
\midrule
1 & Brasil (São Paulo)          & Sí (Zona 1) & 900  \\
2 & Estados Unidos (N.\ Virginia) & Sí (Zona 2) & 1200 \\
3 & Alemania (Frankfurt)        & Sí (Zona 3) & 600  \\
4 & Chile                       & No          & 150  \\
5 & México                      & No          & 400  \\
6 & Argentina                   & No          & 250  \\
7 & Sudáfrica                   & No          & 180  \\
8 & India (Mumbai)              & No          & 700  \\
\bottomrule
\end{tabular}
\captionof{table}{Países atendidos y masa de clientes por país.}
\end{center}

\begin{center}
\renewcommand{\arraystretch}{1.2}
\begin{tabular}{clcccc}
\toprule
\textbf{$j$} & \textbf{Zona AWS} &
\textbf{$S_j$ (alta/serv.)} &
\textbf{$U_j$ (cap./serv.)} &
\textbf{$M_j$ (máx.\ serv.)} &
\textbf{$A_j$ (admin.)} \\
\midrule
1 & Brasil (São Paulo)     & 10 & 300 & 6 & 50 \\
2 & EE.UU.\ (N.\ Virginia) &  8 & 350 & 8 & 40 \\
3 & Alemania (Frankfurt)   & 12 & 300 & 6 & 60 \\
\bottomrule
\end{tabular}
\captionof{table}{Parámetros de cada zona AWS. $S_j$ y $A_j$ en miles de USD;
$U_j$ en miles de clientes.}
\end{center}

\textbf{Compatibilidad por latencia.} Cada país solo puede ser atendido desde
las zonas de su misma región:

\begin{center}
\renewcommand{\arraystretch}{1.2}
\begin{tabular}{lccc}
\toprule
\textbf{País $\backslash$ Zona} &
\textbf{Z1 Brasil} & \textbf{Z2 Virginia} & \textbf{Z3 Frankfurt} \\
\midrule
Brasil          & Sí & Sí & No \\
Estados Unidos  & Sí & Sí & No \\
Alemania        & No & No & Sí \\
Chile           & Sí & Sí & No \\
México          & Sí & Sí & No \\
Argentina       & Sí & Sí & No \\
Sudáfrica       & No & No & Sí \\
India           & No & No & Sí \\
\bottomrule
\end{tabular}
\captionof{table}{Matriz de compatibilidad país--zona (región \emph{Américas}:
Z1, Z2; región \emph{EMEA/Asia}: Z3).}
\end{center}

\textbf{SLA de servicio local:} $\beta = 0{,}5$. Es decir, cada país que posee
zona local (Brasil, EE.UU., Alemania) debe atender al menos el \textbf{50\,\%}
de sus clientes desde su propia zona.

\noindent\rule{\linewidth}{0.4pt}

% =============================================================================
\section*{Parte A \quad Formulación del Modelo \hfill\normalfont\normalsize(35 pts)}
% =============================================================================

Formule un \textbf{modelo de optimización} que minimice el costo total. Su
formulación debe ser completamente algebraica (en notación matemática,
\textbf{no} en código). \textbf{Toda la información necesaria está en el
Contexto y en los Datos del problema}; una parte esencial de esta tarea es que
\textbf{ustedes identifiquen} por su cuenta los componentes del modelo. Su
formulación debe contener:

\begin{enumerate}[label=\alph*), leftmargin=*, topsep=2pt]
    \item \textbf{Conjuntos e índices:} defina los conjuntos que necesite para
          indexar países y zonas (y cualquier subconjunto que le resulte útil).
    \item \textbf{Parámetros:} liste, con su notación, todos los datos del
          problema.
    \item \textbf{Variables de decisión:} identifique y defina \textbf{todas}
          las variables necesarias. Para cada una indique con precisión su
          \textbf{significado}, sus \textbf{unidades}, su \textbf{dominio} y su
          \textbf{tipo} (continua, entera o binaria) según lo que represente.
    \item \textbf{Función objetivo:} plantee la expresión del costo total a
          minimizar a partir de los costos descritos en el Contexto.
    \item \textbf{Restricciones:} formule y \textbf{nombre} \textbf{todas} las
          restricciones necesarias para que el modelo represente fielmente los
          requerimientos de negocio descritos en el Contexto y en los Datos.
          Incluya los dominios de las variables.
\end{enumerate}

\noindent\textit{Sugerencia de método (no de solución): relea con cuidado el
Contexto y la sección de Datos. Cada condición del negocio —cobertura de los
clientes, capacidades, límites, costos que se pagan solo al utilizar un recurso,
latencia y SLA— suele traducirse en una o más restricciones. Decidir cuántas
variables usar, de qué tipo, y qué restricciones escribir es precisamente lo que
se evalúa en esta parte.}

\noindent\rule{\linewidth}{0.4pt}

% =============================================================================
\section*{Parte B \quad Implementación y Resolución en Gurobi \hfill\normalfont\normalsize(35 pts)}
% =============================================================================

Implemente el modelo de la Parte~A en \texttt{gurobipy} y resuélvalo con los
datos entregados. Responda:

\begin{enumerate}[label=\alph*), leftmargin=*, topsep=2pt]
    \item Reporte el \textbf{costo total óptimo} y descompóngalo en costo de
          servidores y costo administrativo.
    \item Reporte, para cada zona: si se \textbf{utiliza}, el \textbf{número de
          servidores} instalados y el \textbf{porcentaje de utilización} de su
          capacidad instalada.
    \item Presente la \textbf{matriz de asignación} (clientes de cada país
          atendidos desde cada zona) en una tabla clara.
    \item Verifique e informe que se cumplen: (i) la cobertura de todos los
          clientes, (ii) el SLA local de Brasil, EE.UU.\ y Alemania, y (iii) la
          compatibilidad por latencia.
    \item Interprete la solución en \textbf{lenguaje de negocios}: ¿qué zonas
          conviene activar y por qué? ¿Alguna zona queda sin utilizar?
\end{enumerate}

\begin{ayuda}{Ayuda -- Sintaxis de \texttt{gurobipy} (referencia del lenguaje)}
\vspace{2pt}
Esta caja es solo una referencia de la \emph{sintaxis} de la herramienta, con
nombres genéricos: \textbf{no} es el modelo, y \textbf{no} indica qué variables
o restricciones usar (eso es parte de lo evaluado). En \texttt{gurobipy} una
variable puede declararse continua, entera o binaria, y las sumas sobre
conjuntos se arman con \texttt{quicksum}:
\begin{lstlisting}
import gurobipy as gp
from gurobipy import GRB, quicksum

m = gp.Model("netflix_aws")

# Ejemplos de sintaxis (nombres genericos, NO son el modelo):
a = m.addVars(Conjunto, lb=0, name="a")             # continua >= 0
b = m.addVars(Conjunto, vtype=GRB.INTEGER, lb=0)    # entera >= 0
c = m.addVars(Conjunto, vtype=GRB.BINARY)           # binaria (0/1)

m.setObjective(quicksum(...), GRB.MINIMIZE)
m.addConstrs((... for k in Conjunto), name="mi_restr")  # nombrela!

m.optimize()
if m.Status == GRB.OPTIMAL:
    print("Optimo:", m.ObjVal)
    print(a[k].X)   # valor de una variable en la solucion
\end{lstlisting}
Construyan el modelo de forma \textbf{algebraica} (con \texttt{addVars} y
\texttt{quicksum} sobre los conjuntos), no escribiendo los números uno a uno
dentro de las restricciones.
\end{ayuda}

\noindent\rule{\linewidth}{0.4pt}
\newpage
% =============================================================================
\section*{Parte C \quad Escenarios y Preguntas Conceptuales \hfill\normalfont\normalsize(30 pts)}
% =============================================================================

% ── C.1 ──────────────────────────────────────────────────────────────────────
\subsection*{C.1 \quad ¿Cuánto pesa el costo administrativo? \hfill\normalfont\normalsize(15 pts)}

\begin{enumerate}[label=\alph*), leftmargin=*, topsep=2pt]
    \item Resuelva nuevamente el modelo \textbf{eliminando el costo
          administrativo} ($A_j = 0$ para toda zona). Compare la solución
          (zonas activas, servidores, costo) con la del caso base.
    \item ¿El costo administrativo cambia \textbf{qué zonas} se activan, o solo
          el valor de la función objetivo? Explique el resultado.
\end{enumerate}

% ── C.2 ──────────────────────────────────────────────────────────────────────
\subsection*{C.2 \quad Preguntas conceptuales \hfill\normalfont\normalsize(15 pts)}

Responda \textbf{razonadamente} (no requieren resolver un nuevo modelo):

\begin{enumerate}[label=\alph*), leftmargin=*, topsep=2pt]
    \item \textbf{Zona vacía.} ¿Bajo qué condiciones una zona termina
          \textbf{sin ningún servidor} en la solución óptima? Relacione su
          respuesta con el costo administrativo $A_j$, la latencia y las
          capacidades disponibles.
    \item \textbf{Servidores ilimitados y sin costo administrativo.} Suponga
          $M_j = \infty$ para toda zona y $A_j = 0$. \textbf{Sin resolver el
          modelo}, ¿cuál sería la estructura de la solución óptima? Indique a
          qué zona convendría asignar cada país y por qué. \emph{Pista: piense
          en el costo por cliente atendido de cada zona, $S_j / U_j$.}
\end{enumerate}

\noindent\rule{\linewidth}{0.4pt}

% =============================================================================
\section*{Entregables}
% =============================================================================

\subsection*{1.\quad Informe (PDF)}

Estructura mínima requerida (cada sección comienza en una nueva página):

\begin{enumerate}[leftmargin=*, topsep=2pt]
    \item \textbf{Portada:} nombre del curso, Tarea~2, integrantes, profesor,
          ayudante y fecha.
    \item \textbf{Introducción:} descripción del problema de negocio y de su
          relevancia.
    \item \textbf{Marco teórico:} conceptos clave utilizados en la formulación
          y resolución (tipos de variables, función objetivo, restricciones,
          factibilidad, etc.).
    \item \textbf{Desarrollo (metodología):} la formulación completa de la
          Parte~A y la explicación de cómo se implementó y resolvió en
          \texttt{gurobipy}.
    \item \textbf{Resultados:} presente \textbf{solo los resultados
          importantes} —solución óptima, tablas de activación y asignación,
          y escenarios de la Parte~C.
          \begin{quote}
          \textit{Ojo:} en esta sección van los resultados puros, sin
          interpretación.
          \end{quote}
    \item \textbf{Análisis de resultados:} interpretación de la solución y de
          los escenarios en lenguaje de negocios.
    \item \textbf{Conclusiones:} sea \textbf{específico}. Evite frases genéricas
          del tipo ``el modelo funciona'' o ``las restricciones son
          importantes''.
    \item \textbf{Referencias} (si aplica).
\end{enumerate}

\begin{ayuda}{Formalidades del informe (se descuentan puntos por incumplimiento)}
\begin{itemize}[topsep=2pt]
    \item Lenguaje técnico y formal, sin errores ortográficos ni gramaticales.
    \item Texto \textbf{justificado}, fuente tamaño \textbf{12}
          (Times New Roman, Arial u otra tipografía estándar).
    \item Cada figura, gráfico o tabla debe llevar \textbf{título y fuente}.
          Si el material es de elaboración propia, debe indicarse
          explícitamente: \emph{Fuente: elaboración propia}. Si proviene de una
          fuente externa, debe citarse.
    \item Sean \textbf{breves y directos}. Se prefiere un análisis corto,
          directo y correcto, por sobre páginas de relleno.
\end{itemize}
\end{ayuda}

\subsection*{2.\quad Cuaderno Jupyter (.ipynb)}

\begin{itemize}[topsep=2pt]
    \item Implementación \textbf{comentada} del modelo en \texttt{gurobipy}
          (Parte~B), con variables enteras y binarias.
    \item Extracción e impresión de la solución: activación, servidores,
          utilización y matriz de asignación.
    \item Código de los escenarios de la Parte~C (comparación con y sin costo
          administrativo).
    \item Código limpio, ordenado y bien estructurado.
\end{itemize}

\medskip
\noindent\textbf{Formato de entrega:} comprima el informe y el notebook en un
archivo \texttt{.zip} o \texttt{.rar} con el nombre \texttt{Tarea2\_GrupoN}.
\textbf{No entregue los archivos por separado.}

\noindent\rule{\linewidth}{0.4pt}

% =============================================================================
\section*{Plazos y Evaluación}
% =============================================================================

\begin{itemize}[topsep=2pt]
    \item \textbf{Plazo de entrega:} Viernes 10 de julio a las 23:59
          (se confirmará en Canvas).
    \item \textbf{Grupos:} 2 a 3 integrantes (ni más ni menos).
    \item \textbf{Inscripción:} anótense en los grupos de Canvas
          (Personas~$\to$ Grupos~$\to$ Tarea~2\,--\,Grupo~n).
\end{itemize}

\vspace{0.5em}

\begin{center}
\renewcommand{\arraystretch}{1.2}
\begin{tabular}{lc}
\toprule
\textbf{Criterio} & \textbf{Ponderación} \\
\midrule
Parte A: Formulación del modelo                      & 35\,\% \\
Parte B: Implementación y resolución en Gurobi       & 35\,\% \\
Parte C: Escenarios y preguntas conceptuales         & 30\,\% \\
\midrule
\textit{Calidad del informe y profundidad del análisis} &
\textit{transversal} \\
\bottomrule
\end{tabular}
\captionof{table}{Criterios de evaluación.}
\end{center}

\vspace{0.3em}
\noindent\textit{La calidad del informe, la profundidad del análisis y la
correcta interpretación de los resultados en lenguaje de negocios impactan
de forma transversal en todas las partes.}

\noindent\rule{\linewidth}{0.4pt}

% =============================================================================
\section*{Notas Importantes}
% =============================================================================

\begin{itemize}[topsep=2pt]
    \item \textbf{Deben usar \texttt{gurobipy}} para la implementación.
    \item Cada variable debe declararse con su \textbf{dominio y tipo
          correctos} (continua, entera o binaria, según lo que represente).
          Identificar los tipos es parte de la evaluación.
    \item \textbf{Nombren} las restricciones al crearlas: facilita la lectura
          de la solución, la depuración y el análisis de infactibilidad.
    \item Construyan el modelo de forma \textbf{algebraica} (con conjuntos,
          \texttt{addVars} y \texttt{quicksum}); no escriban los números uno a
          uno dentro de las restricciones.
    \item Pueden apoyarse en herramientas externas para entender el método,
          pero la \textbf{formulación} y el \textbf{análisis} deben ser propios
          y quedar plasmados en el informe.
    \item En el informe y en el código deben figurar los \textbf{nombres de
          todos los integrantes} del grupo.
\end{itemize}

\vspace{1em}
\begin{center}
\textit{Consultas al correo
\href{mailto:vramireza@udd.cl}{\texttt{vramireza@udd.cl}} o al
\texttt{+56\,9\,9912\,1814}
(respondo siempre que puedo, dentro de horarios razonables).}
\end{center}

% =============================================================================
\end{document}
% =============================================================================
